\section{A shape-sensitive comparison of actual multiplicities}

\subsection*{Finding M-009: bounded weight boxes improve the row gate}
Let \(m\geq2\), \(n>m\), \(d\geq1\), and let
\(\lambda\vdash nd\), with \(r=\ell(\lambda)\leq n^2\).
For a subset \(A\subseteq\{2,\ldots,r\}\), define
\[
 N_A(m,\lambda)=
 \#\{\alpha\in\mathbb Z_{\geq0}^{A}:
       \alpha_i\leq\lambda_i\ (i\in A),\quad |\alpha|\leq m\}.
\]
In particular \(N_\varnothing=1\). Put
\[
 T_m(\lambda)=
 \min_{A\sqcup B=\{2,\ldots,r\}}
       \bigl(N_A(m,\lambda)+N_B(m,\lambda)-1\bigr).
\]
Then
\begin{equation}\label{eq:M009}
 n\geq T_m(\lambda)
 \quad\Longrightarrow\quad
 q^m_\lambda(d[n])\leq K_n(\lambda).
\end{equation}
This compares the actual coordinate-ring multiplicities. It refines
M-008 because deleting the individual bounds \(\alpha_i\leq\lambda_i\)
recovers its two-trie bound.

\begin{proof}
Use \(V=\mathbb C^{n^2}\), with basis \(e_1,\ldots,e_{n^2}\),
and let \(x_i\) denote its dual coordinates. Let \(w\) be a highest-weight
vector of weight \(\lambda\) in \(I(\Omega_n)_d\), for the upper
triangular Borel in \(\mathrm{GL}(V)\). A degree-\(d\) polynomial
function on forms is a polynomial in the coefficient functions
\(a_\eta\), where \(|\eta|=n\); the weight of \(a_\eta\) is
\(\eta\). Every coefficient monomial occurring in \(w\) has total
weight \(\lambda\). All these weights have nonnegative coordinates.
It follows that \(w\) depends only on the first \(r\) coordinates,
and that no coefficient \(a_\eta\) with
\(\eta_i>\lambda_i\) can occur in \(w\).

Restrict a general translate of the padded permanent to those
\(r\) coordinates. It is \(L^{n-m}f\), with \(f\) homogeneous
of degree \(m\). On the nonempty open set on which the coefficient
\(L_1=L(e_1)\) is nonzero, an upper triangular change of basis
aligns \(L\) with \(x_1\). Explicitly, choose an upper unipotent
matrix whose first row is \((1,L_2/L_1,\ldots,L_r/L_1)\);
then \(L\circ g^{-1}=L_1x_1\).
The vector \(w\) is a Borel eigenvector, so this change of basis
rescales its value by a nonzero scalar (it is unchanged for the
unipotent choice). Absorbing \(L_1^{n-m}\) into \(f\), it suffices
to evaluate \(w\) on a form \(x_1^{n-m}f\).

Expand
\[
 f=\sum_{|\alpha|\leq m}
            c_\alpha x_2^{\alpha_2}\cdots x_r^{\alpha_r}
                         x_1^{m-|\alpha|}.
\]
Delete all terms for which some \(\alpha_i>\lambda_i\),
\(i\geq2\), and call the resulting form \(f_\lambda\).
The preceding weight observation proves exactly
\[
       w(x_1^{n-m}f)=w(x_1^{n-m}f_\lambda).
\]
There is no approximation in this equality: every deleted form
coefficient is absent from the polynomial \(w\).

Choose a partition \(A\sqcup B\) attaining \(T_m(\lambda)\).
Make the two trees of M-007, but retain only the multi-indices
bounded by \(\lambda_i\) in their respective coordinates and
of total degree at most \(m\). These vertex sets are closed
under decreasing any coordinate, so the parent rule remains valid.
Their total number of vertices, after identifying their roots, is
\(S=N_A+N_B-1\). Add the coefficient edges for all retained
pairs of total degree at most \(m\), and put \(x_1\) on the
unused-vertex loops. The same cycle-cover argument gives, term by term,
\[
 \det M=\sum_{\substack{|\alpha|\leq m\\
                       \alpha_i\leq\lambda_i\ (i\geq2)}}
                c_\alpha x_2^{\alpha_2}\cdots x_r^{\alpha_r}
                                x_1^{S-|\alpha|}.
\]
Every exponent is nonnegative: the relevant cycle has
\(|\alpha|+1\) distinct vertices in a graph with \(S\) vertices.
This termwise argument is necessary because \(S\) can be smaller
than \(m\). Since \(n\geq S\), adjoining
\(x_1 I_{n-S}\) gives an \(n\times n\) homogeneous linear
determinant equal to \(x_1^{n-m}f_\lambda\).
It belongs to \(\Omega_n\), and \(w\) vanishes on it.

We have proved that \(w\) vanishes on a dense open subset of the
padded-permanent orbit. The orbit is the image of the irreducible
variety \(\mathrm{GL}(V)\), and the condition \(L(e_1)\ne0\)
defines a nonempty open subset of that parameter space. Polynomial
vanishing therefore extends to the entire orbit and its closure.
Hence every highest-weight vector of weight \(\lambda\) in
\(I(\Omega_n)_d\) belongs to \(I(Z_{n,m})_d\).

In the natural decomposition
\(F_d(V)=\bigoplus_\mu S_\mu V\otimes M_\mu\),
the highest-weight space of the \(\lambda\)-isotypic ideal
is its one-dimensional highest-weight line tensored with
\(J_{\Omega_n,\lambda}\). The inclusion just proved gives
\(J_{\Omega_n,\lambda}\subseteq J_{Z_{n,m},\lambda}\).
Subtracting their dimensions from \(\dim M_\lambda\) proves
\eqref{eq:M009} for all copies.
\end{proof}

\paragraph{Explicit thin-tail consequence.}
If \(\lambda_2\leq k\), then every tail row is at most \(k\),
and a balanced split gives
\[
 T_m(\lambda)\leq
  (k+1)^{\lfloor(r-1)/2\rfloor}
       +(k+1)^{\lceil(r-1)/2\rceil}-1.
\]
Thus the right-hand side being at most \(n\) excludes an actual
multiplicity obstruction, even for shapes with more rows than
M-008 covers. The bound counts all points of the box, so discarding
points of total degree greater than \(m\) can improve it further.
For unequal tail rows one may minimize the exact bounded-monomial
counts; no efficient optimization algorithm is claimed.

\paragraph{What this changes and what remains.}
The residual candidate gate now includes \(n<T_m(\lambda)\),
in addition to the independent occurrence, padding, and degree
conditions. Unlike the universal-form method limited by O-010,
this argument only preserves evaluation of a specified highest-weight
space. It does not assert containment of all projected forms in
the determinant closure. It supplies no conclusion for shapes
failing its inequality.

\paragraph{Grothendieck dictionary.}
The coefficient functions are weight sections for the diagonal torus
of \(\mathrm{GL}(V)\); the weight decomposition is the grading/comodule
correspondence of SGA~3, I, Proposition~4.7.3. The upper Borel
acts by its highest-weight character on the line being evaluated.
The deleted coefficients lie outside the finite multigraded support
of that section. The bounded-tree matrices are morphisms of affine
coefficient schemes defined over \(\mathbb Z\), with the displayed
determinant identities integral. The passage from a highest-weight
ideal inclusion to a multiplicity inequality uses semisimplicity
only on the characteristic-zero fiber, as in M-006 and the
Schur-functor dictionary of cote~133, scan~028. No equality of
integral multiplicity modules or of bad-characteristic fibers
is asserted.
